\documentclass[12pt]{article}
\usepackage[utf8]{inputenc}
\usepackage[T1]{fontenc}
\usepackage[margin=1in]{geometry}
\usepackage{setspace}
\usepackage{xcolor}
\usepackage[colorlinks=true, urlcolor=blue!60!black, linkcolor=blue!60!black]{hyperref}
\setlength{\parindent}{2em}
\setlength{\parskip}{1.25em}
\renewcommand{\baselinestretch}{1.0}
\title{\textit{Missouri v. Holland}, 252 U.S. 416 (1920)}
\author{Case Brief}
\date{}

\begin{document}
\maketitle

\section{Case Metadata}
\begin{itemize}
  \item \textbf{Citation:} \textit{State of Missouri v.\ Holland}, 252 U.S.\ 416 (1920) (40 S.\ Ct.\ 382; 64 L.\ Ed.\ 641)
  \item \textbf{Argued / Decided:} Argued March 2, 1920; decided April 19, 1920
  \item \textbf{Vote:} 7--2
  \item \textbf{Author:} Justice Oliver Wendell Holmes Jr.\ (joined by Chief Justice Edward Douglass White, and Justices Joseph McKenna, William R.\ Day, James Clark McReynolds, Louis D.\ Brandeis, and John H.\ Clarke)
  \item \textbf{Concurrences / Dissents:} No concurrences. Justices Willis Van Devanter and Mahlon Pitney dissented without opinion.
  \item \textbf{Chapter / Topic:} Federal Legislative Power --- Treaty Power, Necessary and Proper Clause, and the Tenth Amendment
  \item \textbf{Primary Source URLs:} \href{https://supreme.justia.com/cases/federal/us/252/416/}{Justia U.S.\ Supreme Court Center} $|$ \href{https://www.law.cornell.edu/supremecourt/text/252/416}{Cornell LII}
\end{itemize}

\section{Facts and Procedural History}
In 1913, Congress enacted the Migratory Bird Act (Weeks-McLean Law), attempting unilaterally---without any underlying treaty---to regulate the hunting of migratory birds across the United States. Two federal district courts struck down the 1913 Act as exceeding Congress's enumerated powers and violating the Tenth Amendment (\textit{United States v.\ Shauver}, 214 F.\ 154 (E.D.\ Ark.\ 1914); \textit{United States v.\ McCullagh}, 221 F.\ 288 (D.\ Kan.\ 1915)), reasoning under \textit{Geer v.\ Connecticut}, 161 U.S.\ 519 (1896), that wild game belonged to the states in their sovereign capacity for the benefit of their people.

While the constitutional challenge to the 1913 statute was pending, the United States and Great Britain (acting on behalf of Canada) concluded the Migratory Bird Treaty of August 16, 1916 (proclaimed December 8, 1916), 39 Stat.\ 1702. The treaty recited that many species of birds in their annual migrations traversed parts of the United States and Canada, were of great value as a source of food and in destroying insects injurious to vegetation, and were in danger of extermination through lack of adequate protection. It provided for specified closed hunting seasons and bound both nations to take or propose to their lawmaking bodies the necessary measures to carry the treaty out.

To give domestic effect to the treaty, Congress enacted the Migratory Bird Treaty Act of July 3, 1918, ch.\ 128, 40 Stat.\ 755, prohibiting the killing, capturing, or selling of covered migratory birds except as permitted by regulations made by the Secretary of Agriculture (proclaimed July 31 and October 25, 1918). The State of Missouri filed a bill in equity in the U.S.\ District Court for the Western District of Missouri to enjoin Ray P.\ Holland, a U.S.\ Game Warden, from enforcing the 1918 Act and its regulations. Missouri contended that the statute unconstitutionally interfered with powers reserved to the states by the Tenth Amendment and invaded the state's sovereign and pecuniary title to wild birds within its borders. The District Court sustained the government's motion to dismiss on the ground that the Act was constitutional (258 F.\ 479 (W.D.\ Mo.\ 1919)), and Missouri appealed directly to the U.S.\ Supreme Court.

\section{Issue(s) Presented}
\begin{enumerate}
  \item Does a state have standing in equity to enjoin federal enforcement of a statute allegedly infringing its quasi-sovereign authority to regulate wild game?
  \item May Congress, acting under the Necessary and Proper Clause (Article~I, Section~8, Clause~18) to implement a valid treaty made under Article~II, Section~2, enact domestic legislation regulating a subject that might otherwise fall outside Congress's standalone Article~I enumerated powers and within the Tenth Amendment's reserved powers of the states?
\end{enumerate}

\section{Holding and Rule}
Yes to both. Missouri's assertion of its quasi-sovereign right to regulate game sufficed to invoke equity jurisdiction apart from any pecuniary interest. On the merits, the Migratory Bird Treaty was a valid exercise of the federal treaty-making power under Article~II, Section~2, and the Migratory Bird Treaty Act was valid under Article~I, Section~8, as a necessary and proper means of executing that treaty. \textbf{Rule:} Because the treaty-making power is expressly delegated to the United States and treaties are the supreme law of the land under Article~VI, a valid treaty and its implementing legislation are not constrained by what an unaided Act of Congress could accomplish under Article~I and do not violate the Tenth Amendment. Decree affirmed.

\section{Reasoning}
Writing for a seven-Justice majority in one of the most concise and consequential opinions in constitutional law, Justice Holmes dismantled Missouri's Tenth Amendment challenge through four interlocking steps:

\begin{itemize}
  \item \textbf{Express Delegation and Necessary-and-Proper Implementation:} Holmes began by observing that the Tenth Amendment reserves only powers ``not delegated to the United States.'' Because Article~II, Section~2 expressly delegates the power to make treaties, and Article~VI declares treaties to be the supreme law of the land, invoking the Tenth Amendment does not answer the question. In a pivotal single-sentence bridge, Holmes declared: ``If the treaty is valid, there can be no dispute about the validity of the statute under Article~I, \S~8, as a necessary and proper means to execute the powers of the Government.''
  \item \textbf{Supremacy Clause Textual Asymmetry:} Holmes underscored a textual distinction in Article~VI: ``Acts of Congress are the supreme law of the land only when made in pursuance of the Constitution, while treaties are declared to be so when made under the authority of the United States.'' While noting it was ``open to question'' whether ``under the authority of the United States'' meant more than the formal constitutional steps required to make a treaty, Holmes disclaimed any suggestion that the treaty power is wholly unlimited---insisting only that its limits must be ascertained ``in a different way'' than by testing it against unaided Article~I legislation.
  \item \textbf{Constitutional Evolution and Collective-Action Necessity:} In a classic articulation of living constitutionalism, Holmes wrote that the Constitution is ``a constituent act'' that ``called into life a being the development of which could not have been foreseen completely by the most gifted of its begetters.'' Constitutional boundaries must therefore be evaluated ``in the light of our whole experience, and not merely in that of what was said a hundred years ago.'' Because the treaty contravened no express ``prohibitory words'' in the Constitution, the sole question was whether it was barred by ``some invisible radiation from the general terms of the Tenth Amendment.'' Where a problem of ``the sharpest exigency for the national wellbeing'' cannot be solved by individual states acting alone, a civilized national government must possess the power to act in concert with another sovereign.
  \item \textbf{Rejection of State ``Title'' to Migratory Wildlife:} Holmes dismissed Missouri's claim of proprietary ownership over wild birds as leaning ``upon a slender reed.'' Possession is the beginning of ownership, and migratory birds are transitory creatures ``that yesterday had not arrived, tomorrow may be in another State, and, in a week, a thousand miles away.'' Citing early precedents in which treaties validly overrode state statutes of limitations, confiscation laws, and escheat of land (\textit{Ware v.\ Hylton}, 3 U.S.\ (3 Dall.) 199 (1796); \textit{Hopkirk v.\ Bell}, 7 U.S.\ (3 Cranch) 454 (1806); \textit{Chirac v.\ Chirac}, 15 U.S.\ (2 Wheat.) 259 (1817); \textit{Geofroy v.\ Riggs}, 133 U.S.\ 258 (1890)), Holmes concluded that protecting a vital food supply and agricultural resource required international cooperation; relying on individual states was ``vain.''
\end{itemize}

\section{Separate Opinions (Concurrences and Dissents)}
No concurring opinions were filed. Justices Willis Van Devanter and Mahlon Pitney dissented without opinion. Although neither Justice published reasons in 1920, subsequent scholarship and the three concurring opinions in \href{https://supreme.justia.com/cases/federal/us/572/844/}{\textit{Bond v.\ United States}}, 572 U.S.\ 844 (2014), articulated the structural objections underlying their silent dissent: namely, that allowing a non-self-executing treaty to expand Congress's domestic legislative authority creates a loophole around enumerated Article~I powers and the Tenth Amendment.

\section{Key Quotations}
\begin{enumerate}
  \item ``If the treaty is valid, there can be no dispute about the validity of the statute under Article~I, \S~8, as a necessary and proper means to execute the powers of the Government.'' --- Justice Holmes, 252 U.S.\ at 432 (\href{https://supreme.justia.com/cases/federal/us/252/416/}{Justia}).
  \item ``Acts of Congress are the supreme law of the land only when made in pursuance of the Constitution, while treaties are declared to be so when made under the authority of the United States.'' --- Justice Holmes, 252 U.S.\ at 433 (\href{https://supreme.justia.com/cases/federal/us/252/416/}{Justia}).
  \item ``It was enough for them to realize or to hope that they had created an organism; it has taken a century and has cost their successors much sweat and blood to prove that they created a nation. The case before us must be considered in the light of our whole experience, and not merely in that of what was said a hundred years ago.'' --- Justice Holmes, 252 U.S.\ at 433 (\href{https://supreme.justia.com/cases/federal/us/252/416/}{Justia}).
  \item ``The treaty in question does not contravene any prohibitory words to be found in the Constitution. The only question is whether it is forbidden by some invisible radiation from the general terms of the Tenth Amendment.'' --- Justice Holmes, 252 U.S.\ at 433--434 (\href{https://supreme.justia.com/cases/federal/us/252/416/}{Justia}).
  \item ``To put the claim of the State upon title is to lean upon a slender reed. Wild birds are not in the possession of anyone, and possession is the beginning of ownership. The whole foundation of the State's rights is the presence within their jurisdiction of birds that yesterday had not arrived, tomorrow may be in another State, and, in a week, a thousand miles away.'' --- Justice Holmes, 252 U.S.\ at 434 (\href{https://supreme.justia.com/cases/federal/us/252/416/}{Justia}).
\end{enumerate}

\section{Doctrinal Significance and Modern Legacy}
\textit{Missouri v.\ Holland} remains the foundational precedent governing the intersection of the Treaty Power, the Necessary and Proper Clause, and American federalism, though its practical reach has been substantially cabined over the past century.

\begin{itemize}
  \item \textbf{Individual-Rights Limits and the Bricker Amendment:} Holmes's observation about Article~VI's ``under the authority of the United States'' wording fueled mid-twentieth-century fears that international treaties could override domestic constitutional rights or state sovereignty, sparking the proposed (but narrowly defeated) \href{https://www.law.cornell.edu/constitution-conan/article-2/section-2/clause-2/congressional-implementation-of-treaties}{Bricker Amendment} in the 1950s. The Court put the individual-rights fear to rest in \href{https://supreme.justia.com/cases/federal/us/354/1/}{\textit{Reid v.\ Covert}}, 354 U.S.\ 1, 16--18 (1957), where Justice Black's plurality held that ``no agreement with a foreign nation can confer power on the Congress, or on any other branch of Government, which is free from the restraints of the Constitution'' (barring peacetime courts-martial of civilian military dependents overseas), while distinguishing \textit{Holland} on the ground that the Tenth Amendment does not bar a delegated treaty power.
  \item \textbf{Non-Self-Execution Doctrine:} In \href{https://supreme.justia.com/cases/federal/us/552/491/}{\textit{Medell\'{i}n v.\ Texas}}, 552 U.S.\ 491 (2008), the Court reinforced that non-self-executing treaties do not create domestically enforceable federal law absent implementing legislation by Congress---heightening the practical stakes of \textit{Holland}'s Necessary and Proper Clause holding whenever Congress legislates to execute modern multilateral agreements.
  \item \textbf{The \textit{Bond v.\ United States} Challenge (2011 \& 2014):} \textit{Holland} faced its most direct modern challenge in the \textit{Bond} litigation. After \href{https://supreme.justia.com/cases/federal/us/564/211/}{\textit{Bond v.\ United States} (\textit{Bond I})}, 564 U.S.\ 211 (2011), unanimously held that an individual criminal defendant has standing to raise a Tenth Amendment challenge, the Court reached the merits in \href{https://supreme.justia.com/cases/federal/us/572/844/}{\textit{Bond v.\ United States} (\textit{Bond II})}, 572 U.S.\ 844 (2014). Carol Anne Bond, a Pennsylvania microbiologist who spread toxic chemicals on the car door, mailbox, and doorknob of her husband's pregnant lover (causing a minor thumb burn), was prosecuted under the federal Chemical Weapons Convention Implementation Act of 1998 (18 U.S.C.\ \S~229). Because the government waived any Commerce Clause defense, the case squarely presented whether \textit{Holland} allowed a treaty-implementing statute to federalize a purely local assault.
  \begin{itemize}
    \item \textbf{Chief Justice Roberts's Majority (6--3 on rationale; 9--0 on judgment):} Invoking constitutional avoidance and federalism clear-statement canons (\textit{Gregory v.\ Ashcroft}; \textit{United States v.\ Bass}), Roberts declined to revisit or overrule \textit{Holland}. Instead, he construed \S~229 narrowly as not reaching ``purely local crimes'' absent a clear indication of congressional intent, observing that the ``global need to prevent chemical warfare does not require the Federal Government to reach into the kitchen cupboard.''
    \item \textbf{The Three-Justice Originalist Critique:} Concurring only in the judgment, Justice Scalia (joined by Justice Thomas, and by Justice Alito as to Part~I) attacked \textit{Holland}'s Necessary and Proper Clause bridge as an ``unreasoned and citation-less'' \textit{ipse dixit}. Textually, Scalia argued that Article~I, \S~8, cl.\ 18 authorizes laws necessary and proper for carrying into execution the power ``\textit{to make} Treaties'' (e.g., funding negotiators), \textit{not} for implementing treaties already made; once a treaty is ratified, Congress must rely on its independent Article~I powers. Otherwise, Scalia warned, the President and Senate could overturn \href{https://supreme.justia.com/cases/federal/us/514/549/}{\textit{United States v.\ Lopez}}, 514 U.S.\ 549 (1995), simply by negotiating a gun-free-school-zones treaty with Latvia. Justice Thomas (joined by Scalia, and in Parts~I--III by Alito) wrote separately to argue from Founding-era history that the Treaty Power itself is substantively confined to ``matters of international intercourse'' rather than purely domestic affairs---though Thomas notably observed that \textit{Holland} itself was consistent with that limit because migratory birds were ``creatures in international transit.'' Justice Alito agreed that the Treaty Power reaches only ``matters of legitimate international concern.''
  \end{itemize}
  \item \textbf{Status Through 2026:} As of 2026, \textit{Missouri v.\ Holland} remains formally good law---neither overruled nor formally narrowed on its constitutional holding. Yet its modern operation is doubly constrained: (1) under \textit{Bond II}, courts vigorously deploy federalism clear-statement canons to narrow treaty-implementing statutes that intrude on traditional state police powers; and (2) given a contemporary Supreme Court majority deeply receptive to originalist and structural-federalism limits (cf.\ the Tenth Amendment and enumerated-powers debates in \href{https://supreme.justia.com/cases/federal/us/599/21-376/}{\textit{Haaland v.\ Brackeen}}, 599 U.S.\ 255 (2023)), a direct constitutional challenge to an unambiguous statute implementing a multilateral treaty on a purely domestic subject would place \textit{Holland}'s survival in serious jeopardy.
\end{itemize}

\section{Pedagogical Critiques and Exam Hypotheticals}
\textbf{Pedagogical Critiques:}
\begin{itemize}
  \item \textbf{Textual Slippage in Article~VI and Article~I:} Scholars (notably Nicholas Quinn Rosenkranz, cited by Justice Scalia in \textit{Bond II}) highlight two textual vulnerabilities in Holmes's opinion. First, the Supremacy Clause's phrase ``under the Authority of the United States'' (as opposed to ``in Pursuance thereof'') was historically drafted to preserve pre-1787 treaties made under the Articles of Confederation---not to exempt future treaties from constitutional limits. Second, Holmes never justified why the Necessary and Proper Clause empowers Congress to execute a treaty's \textit{substantive obligations} rather than merely assisting the Article~II process of \textit{making} treaties.
  \item \textbf{Functionalism Without a Stopping Point:} Holmes's collective-action rationale is compelling for migratory birds crossing the U.S.--Canada border, where no single state can prevent extinction. Yet his opinion supplies no judicially manageable standard for distinguishing genuine international collective-action problems from pretextual treaties negotiated to bypass Article~I limits.
\end{itemize}

\textbf{Exam Hypothetical:} Congress previously enacted a federal statute banning corporal punishment in all public and private K--12 schools, which a federal court invalidated under \textit{United States v.\ Lopez} as regulating noneconomic local conduct. The President subsequently negotiates, and two-thirds of the Senate ratifies, the ``International Convention on the Rights of the Child,'' which obligates signatory nations to enact domestic laws banning corporal punishment in all schools. Congress then reenacts the identical school-discipline ban, expressly stating that it applies to every local public and private school nationwide and citing the treaty as its constitutional authority. A private school in Texas challenges the statute.
\begin{itemize}
  \item \textbf{Government's Best Argument (Application of \textit{Missouri v.\ Holland}):} Under \textit{Holland}, Congress's power to enact legislation necessary and proper to execute a valid treaty is not limited by what an unaided Act of Congress could accomplish under the Commerce Clause. Because Congress included an unmistakable clear statement reaching local schools, \textit{Bond II}'s statutory-avoidance canon cannot defeat the statute. Furthermore, the statute violates no affirmative Bill of Rights prohibition (\textit{Reid v.\ Covert}), and international human-rights norms protecting children are a well-established subject of multilateral diplomacy.
  \item \textbf{Challenger's Best Argument (Opposite-Side Contrast via \textit{Bond II} Concurrences):} Because Congress's explicit text forecloses \textit{Bond II}'s avoidance path, the Court must confront the constitutional questions left open in \textit{Bond II}. First, under Justice Thomas's and Justice Alito's \textit{Bond II} framework, the Treaty Power is originally limited to ``matters of international intercourse'' or ``legitimate international concern.'' Unlike migratory birds physically crossing international borders in \textit{Holland}, intrastate school discipline involves no cross-border transit or genuine international collective-action failure; upholding it would allow the President and Senate to overturn \textit{Lopez} at will. Second, under Justice Scalia's textualist reading of the Necessary and Proper Clause, Congress may assist the \textit{making} of treaties, not bootstrap a non-self-executing treaty into a freestanding domestic police power beyond Article~I, \S~8.
\end{itemize}

\end{document}
